\clearpage
\begin{summarybox}
	\textbf{\textcolor{CSummary}{一句话核心}}

	\textbf{双曲线切线方程可由原方程通过“半代换”得到：将 $x^{2}$ 换为 $x_{0}x$，$y^{2}$ 换为 $y_{0}y$，方程结构不变。}

	\medskip
	\textbf{\textcolor{CSummary}{使用条件}}
	\begin{itemize}[leftmargin=*, itemsep=0.5em, topsep=0.4em]
		\item \textbf{条件 1（切点在线上）：} 点 $P(x_0,y_0)$ 必须位于双曲线上。
		\item \textbf{条件 2（标准形式）：} 双曲线方程必须为 $\frac{x^{2}}{a^{2}}-\frac{y^{2}}{b^{2}}=1$ 或 $\frac{y^{2}}{a^{2}}-\frac{x^{2}}{b^{2}}=1$ 的标准形式。
	\end{itemize}

	\medskip
    \textbf{\textcolor{CSummary}{三种证明路线（固定顺序）}}
    \begin{enumerate}[leftmargin=*, itemsep=0.4em, topsep=0.3em]
      \item \textbf{最低门槛｜平方差 $+$ 割线逼近：}两点代入后相减，
      得 $k_{PQ}=\dfrac{b^2(x+x_0)}{a^2(y+y_0)}$；让 $Q$ 靠近 $P$，
      得 $k=\dfrac{b^2x_0}{a^2y_0}$（$y_0\ne0$）。
      \item \textbf{最快｜导数 $+$ 点斜式：}直接求 $y'=\dfrac{b^2x}{a^2y}$，
      再由点斜式和切点在曲线上得到切点式。
      \item \textbf{常规｜联立 $+$ 判别式：}设 $y=kx+t$，
      联立为二次方程并令 $\Delta=0$，解得同样的 $k$。
    \end{enumerate}
    \textbf{三法共同补充：}若 $y_0=0$，则 $x_0=\pm a$，切线是 $x=x_0$。

    \medskip
    \textbf{\textcolor{CSummary}{关键提醒}}
	\begin{itemize}[leftmargin=*, itemsep=0.5em, topsep=0.4em]
		\item \textbf{易错点（焦点混淆）：} 注意双曲线开口方向，对应代换公式中 $x$ 和 $y$ 的位置。
		\item \textbf{判别式前提：} 联立后必须确实得到一元二次方程；二次项消失时即使只有一个交点，也不能直接说相切。
		\item \textbf{检查项（竖直切线）：} 当 $y_0=0$ 时切线斜率不存在，但公式依然适用，用导数法时须单独讨论，不能把零代入斜率式的分母。
	\end{itemize}
\end{summarybox}
